\startsectionlevel[title={{\bf 94. くらい VS ほど}},reference={くらい-vs-ほど}]

\goto{{\bf くらい (kurai) VS ほど (hodo) What they mean and why they mean it.}}[url(https://www.youtube.com/watch?v=6PEQTcDnbBk&ab_channel=OrganicJapanesewithCureDolly)]

こんにちは。

Today we're going to talk about two words that often cause people difficulty and confusion in immersion because they're often not explained very clearly, so it's difficult sometimes to understand what exactly it is that they're doing.

{\bf The two words are くらい and ほど.}

{\externalfigure[../media/image1072.webp]}

Now, I've had a lot of requests for くらい, here and on my Patreon, and quite a few requests for ほど. And I'm going to deal with them together because I think they can throw a little light on each other. And before we begin, I'll just say that if you have any words that are causing you trouble, bringing you to a halt in understanding things in your immersion, please put them in the Comments below and I'll either reply right away or add them to my list of things to deal with.

\thinrule

Right, so, ほど and くらい. {\bf They're both nouns, of course.}

{\bf Practically everything in Japanese that isn't a verb or an adjective, as we know, is in practice a noun or a noun-like entity.}

So, what is a ほど? What is a くらい?

{\bf Both of them represent an extent or an amount of\ldots{} anything: time, money, distance, etc.}

{\externalfigure[../media/image343.webp]}

{\bf However, the main difference between them is that ほど tends to imply that the amount or degree involved is great,}

\thinrule

{\bf and くらい either marks it as indefinite, as something we're not sure about, or in some cases which we'll talk about a bit later, that it is small, that it's less than we might have expected or hoped.}

So, to illustrate this difference, let's look at them both in a very simple construction. {\bf これくらい means: about this much, about this amount, about this far, about this fast,\crlf
about this temperature, whatever.}

{\bf これほど means: this much, this amount, this hot, this cold, but it's stressing that this is a large extent, a large degree.}

\startsectionlevel[title={ほど},reference={ほど}]

So, if we say {\bf これほど悪い}とは思わなかった, we're saying, I didn't think it was {\bf this bad}.

{\bf So, the これ (this) modifying ほど is not only saying this bad, it's saying that this bad is a lot, this is more than we might have expected or desired.}

{\externalfigure[../media/image814.webp]}

{\bf Negatively}, if we say {\bf それほど}悪くない, we're saying it's not {\bf that} bad.

{\externalfigure[../media/image55.webp]}

Now, {\bf what's それ referring to in a case like this?}

{\bf It can be referring to something we've already referred to, to a degree or an extent that we've already been discussing,} but also, just as in English, {\bf それ (or that) can simply mean an amount that we might have expected.}

\thinrule

{\bf So, again, the thing modifying ほど, the それ, represents a great amount, an upper limit, and then we're falling short of that upper limit marked by ほど.}

So, {\bf それほど}悪くない means not {\bf that} bad --- {\bf either not that bad} as we have been discussing or as is implied in the conversation, {\bf or simply not that bad}, which works exactly the same as it does in English --- {\bf the that when there's no reference point simply means what we might have expected or what it might have been.}

So if we say in English it's not that bad, we mean {\bf it's not very bad, it's not as bad as it might have been.}

\thinrule

And the same with {\bf それほど}悪くない. {\bf If we're not referring to a specific, defined that, then we're saying it's not }very** bad / it's not {\bf all that} bad.**

And again, {\bf ほど is often used to mark one side of a comparison, and once again ほど represents the greater side, the upper limit.}

So if we say 私は彼女{\bf ほど}若くない we're saying I'm not {\bf as} young {\bf as} she is,

{\externalfigure[../media/image196.webp]}

I'm not young {\bf to the extent that} she is, and {\bf her extent} is {\bf the higher extent} {\bf and therefore it's the one modifying ほど.}

\startsectionlevel[title={ほど as a noun},reference={ほど-as-a-noun}]

Now, {\bf ほど can be used as an almost independent noun representing the upper limit of something.}

So if we say 不可能{\bf なほど}, we're saying (不可能: 可能 means possible; 不可能 means impossible) 不可能な{\bf ほど} means {\bf the upper limit of} impossibility.

{\bf This is not just 不可能, not just impossible, it's the ほど of 不可能, it's the upper limit of 不可能, it's the height of impossibility.}

{\externalfigure[../media/image1043.webp]}

Or 不快{\bf なほど}失礼 --- 不快 is unpleasant, dislikeable; 失礼, of course, is rudeness. This is rudeness {\bf to the upper limit of} unpleasantness / dislikeableness.

\thinrule

And these {\bf ほどs where we use an adjectival noun and add なほど} --- {\bf the upper limit of that adjectival noun, using that adjectival noun to modify the ほど} --- {\bf very often they tend to be negative things but they can be more neutral,} for example, we might say 不思議{\bf なほど} --- {\bf the upper limit of} mysteriousness, absolutely mysterious.

This phrase above is apparently an expression that can also mean wondrous/marvellous

{\bf And another commonly used expression with ほど is ほどがある, and this really means an upper limit exists.}

{\bf And this is usually used in complaining about something.}

So we say 冗談にも{\bf ほどがある}, which is literally saying even a joke {\bf has an upper limit}, and what that usually means is, what you're saying / what you're doing {\bf goes beyond} a joke.

{\externalfigure[../media/image1028.webp]}

{\externalfigure[../media/image549.webp]}

Jokes have {\bf an upper limit} and you've passed it. バカにも{\bf ほどがある} (stupidity {\bf has an upper limit}), which is a bit like saying, Even you can't be that stupid or you've passed {\bf the upper limit of} stupidity.

\stopsectionlevel

\startsectionlevel[title={ほど for forming comparisons},reference={ほど-for-forming-comparisons}]

{\bf And perhaps the commonest use of ほど,} the one that you'll probably see all the time, {\bf is its use in forming comparisons.}

Now, {\bf there are various kinds of comparison in Japanese, but ほど is used particularly when we're stressing the extreme quality of something.}

{\bf It's often used for comparisons that are essentially exaggerations but making the point that something is very much whatever it is we're saying it is.}

So, we might say, 死ぬ{\bf ほど}暑い (it's hot {\bf to the extent of} dying); 肌を刺す{\bf ほど}寒い (it's cold {\bf to the extent of} penetrating the skin);

{\externalfigure[../media/image933.webp]}

信じられない{\bf ほど}美しい景色 (scenery that's beautiful {\bf to the extent of} being unbelievable).

\thinrule

{\bf And this ほど which is used for exaggeration, for literary purposes, just for making a point very strongly, is probably the ほど that you're going to see more of the time than any other.}

{\bf It's a very common use of ほど.}

\stopsectionlevel

\stopsectionlevel

\startsectionlevel[title={くらい},reference={くらい}]

So, let's look at くらい. {\bf くらい is most commonly used for approximation.}

So we might say, 到着するのは八時{\bf ぐらい}です --- we're saying we're arriving {\bf at approximately} eight o'clock.

{\externalfigure[../media/image1042.webp]}

八百円{\bf ぐらい}です (it's {\bf around} eight hundred yen). And this is very straightforward.

くらい as you can see can sometimes change in sound to ぐらい in some phrases.

{\em As to when, that is apparently no longer as clear, as can be read \goto{{\bf here}}[url(https://japanese.stackexchange.com/questions/433/is-there-a-rule-for-when-to-use-\%E3\%81\%8F\%E3\%82\%89\%E3\%81\%84-vs-\%E3\%81\%90\%E3\%82\%89\%E3\%81\%84)]. But wouldn't worry much.}

{\externalfigure[../media/image938.webp]}

{\bf We use it whenever we want to talk about an amount, a degree, an extent, and make it vague, make it clear that we're not talking about an exact amount or an exact degree.}

\thinrule

{\bf And the only thing to bear in mind here is that sometimes Japanese people will use this when in fact the thing in question is exact.}

So, for example, we might say, どの{\bf くらい}掛かりますか? ({\bf about} how much will it cost?) {\bf And we may be asking for an estimate, if it's the sort of thing where you ask for an estimate, but we may also be asking the price in a case where we can expect the seller to know the exact price, but somehow it seems a bit less pressing to ask for an approximate price rather than the exact one.}

\thinrule

{\bf And people will also sometimes do this when they're talking about a time of arrival or\ldots{} anything.}

{\bf And this is partly because Japanese people are on the whole very precise, so if they think there might be any chance of the time being slightly different or any fact being slightly different from what they're saying, it feels a bit safer to vague it up a little bit so that you haven't committed yourself to something exact that might not turn out to be exactly true.}

\startsectionlevel[title={くらい as a small extent of something = at least},reference={くらい-as-a-small-extent-of-something-at-least}]

Now, {\bf the other use of くらい (or ぐらい)}, as I mentioned before, {\bf talks about an extent and makes the point that the extent is what we consider to be a small one rather than the large one that ほど implies.}

So, if we say 今日{\bf くらい}家族と過ごそう, we're saying Today, {\bf at least}, let's spend as a family.

{\externalfigure[../media/image842.webp]}

{\bf And it's that くらい which gives us the at least meaning.}

{\bf It shows that when we are saying today we're not just saying today, we're implying that maybe we don't spend much time as a family, maybe we won't be able to spend time as a family again, but today at least let's spend time as a family.}

\startsectionlevel[title={くらい as complaining about something},reference={くらい-as-complaining-about-something}]

{\bf And very often this くらい is used when one is complaining about something and when one is really saying that not only is the くらい small, but even that small くらい wasn't forthcoming.}

{\bf And when it's used like that it's often accompanied by another at least expression like 少なくとも or せめて, and often marked at the end by a marker of regret such as のに.}

\thinrule

So, we might say, 少なくとも「ありがとう」{\bf ぐらい}言ってくれてもいいのに, which literally is at least {\bf as little as} thank you if you kindly said would be good but\ldots{}

{\externalfigure[../media/image799.webp]}

In natural English, we'd probably say You might at least say thank you.

What we're here saying is that {\bf as little as} thank you, if you had kindly said, would be good, {\bf and we're adding のに, which,} as \goto{{\bf I've explained elsewhere}}[url(https://www.youtube.com/watch?v=Au5JOtcwE7A)], {\bf is in fact a contrastive clause connector.}

{\bf It's used to connect one logical clause to another to make a complete sentence with the implication of but, that the second clause somewhat contrasts or contradicts with the first, but it's often used like this as a sentence ender.}

And the implication here is, it would be good if you said {\bf as little as} thank you, {\bf but} (you don't even do that).

So, I hope this clarifies くらい and ほど.

\stopsectionlevel

\stopsectionlevel

\stopsectionlevel

\stopsectionlevel
